\RequirePackage{iftex}
\ifXeTeX
\else
    \GenericError{}{编译错误：本模板必须使用 XeLaTeX 编译！}
    {请将编译器切换为 XeLaTeX 重新编译。}{}
    \stop
\fi

\documentclass[openany]{book}

% ==================== 基础宏包 ====================
\usepackage[UTF8]{ctex}
\usepackage{graphicx,xcolor,tocloft}
\usepackage{fancyhdr,geometry,titlesec}
\usepackage{amsmath,amssymb}
\usepackage{tikz}
\usepackage{amsthm}
\usepackage{enumitem}
\usepackage{silence}
\usepackage{array}
\usepackage{booktabs}
\usepackage{hyperref}
\hypersetup{
    colorlinks=true,
    urlcolor=myblue,
    linkcolor=myblue
}

% ==================== 页面边距与颜色 ====================
\geometry{
    left=2cm,
    right=2cm,
    top=2cm,
    bottom=2cm
}

\definecolor{myblue}{RGB}{0,83,159}

% ==================== 字体警告处理 ====================
\WarningFilter{latexfont}
{Font shape `TU/FandolSong-Regular(0)/b/it' undefined}

% ==================== 定理环境设置 ====================
\newtheoremstyle{mydefinition}
  {\topsep}
  {\topsep}
  {}
  {0pt}
  {\heiti\bfseries}
  {.}
  {.5em}
  {}

\newtheoremstyle{myplain}
  {\topsep}
  {\topsep}
  {}
  {0pt}
  {\heiti\bfseries}
  {.}
  {.5em}
  {}

\theoremstyle{mydefinition}
\newtheorem{definition}{定义}
\newtheorem{example}{例}

\theoremstyle{myplain}
\newtheorem{theorem}{定理}
\newtheorem{lemma}[theorem]{引理}
\newtheorem{corollary}[theorem]{推论}
\newtheorem{proposition}[theorem]{命题}

% ==================== 目录设置 ====================
\renewcommand{\cfttoctitlefont}
{\hfill\color{myblue}\heiti\Large}

\renewcommand{\cftaftertoctitle}{\hfill}

\renewcommand{\cftpartpagefont}{}
\renewcommand{\cftpartafterpnum}{\hspace{0pt}}
\renewcommand{\cftchapfont}{\normalfont}
\renewcommand{\cftsecfont}{\normalfont}
\renewcommand{\cftsubsecfont}{\normalfont}

\setlength{\cftbeforetoctitleskip}{-0.5em}
\setlength{\cftaftertoctitleskip}{0.5em}
\setlength{\cftbeforechapskip}{3pt}
\setlength{\cftbeforesecskip}{0pt}
\setlength{\cftbeforesubsecskip}{0pt}

% ==================== 页眉页脚 ====================
\setlength{\headheight}{15pt}

\pagestyle{fancy}
\fancyhf{}

\fancyhead[LE,RO]
{\color{myblue}\leftmark}

\fancyfoot[LE,RO]
{\color{myblue}\thepage}

\renewcommand{\headrule}
{{\color{myblue}\hrule height 0.4pt width\headwidth}}

\renewcommand{\footrule}
{{\color{myblue}\hrule height 0.4pt width\headwidth}}

\fancypagestyle{plain}{%
    \fancyhf{}%
    \fancyhead[LE,RO]
    {\color{myblue}\leftmark}%
    \fancyfoot[LE,RO]
    {\color{myblue}\thepage}%
    \renewcommand{\headrule}
    {{\color{myblue}\hrule height 0.4pt width\headwidth}}%
    \renewcommand{\footrule}
    {{\color{myblue}\hrule height 0.4pt width\headwidth}}%
}

% ==================== 总标题格式 ====================
\titleformat{name=\chapter,numberless}
    {\color{myblue}\heiti\Large\centering}
    {}
    {0pt}
    {}

\titlespacing*{\chapter}
{0pt}{-4em}{1em}

% ==================== 一级标题格式 ====================
% 四个主要部分依次编号
\renewcommand{\thesection}{\arabic{section}}

\titleformat{\section}
    {\color{myblue}\heiti\large}
    {\thesection}
    {0.5em}
    {}

\titlespacing*{\section}
{0pt}{0.8em}{0.5em}

% ==================== 二级标题格式 ====================
% 编号形式为 3.1、3.2、4.1 等
\renewcommand{\thesubsection}
{\thesection.\arabic{subsection}}

\titleformat{\subsection}
    {\color{myblue}\heiti\normalsize}
    {\thesubsection}
    {0.5em}
    {}

\titlespacing*{\subsection}
{0pt}{0.7em}{0.5em}

% ==================== 图题设置 ====================
\renewcommand{\figurename}{图}

% ==================== 表格格式 ====================
\renewcommand{\arraystretch}{1.4}

% ======================================================
%                       正文开始
% ======================================================

% ==================== 本文排版设置 ====================
\setlength{\emergencystretch}{2em}
\setlist[itemize]{leftmargin=2em,itemsep=0.2em,topsep=0.3em}
\newcommand{\Sother}{\mathbf S_{\mathrm{other}}}
\newcommand{\fixed}{\text{不变}}

\begin{document}
\raggedbottom

% ==================== 封面 ====================
\begin{titlepage}
    \thispagestyle{empty}
    \centering
    \vspace*{2cm}
    {\color{myblue}\rule{\textwidth}{2pt}}\par
    \vspace{1.5cm}
    {\color{myblue}\Huge\textbf{Four-State Method for Magnetic Interactions}}\par
    \vspace{0.5cm}
    \vspace{0.5cm}
    {\color{myblue}\Large 对称交换、DM 相互作用与单离子各向异性}\par
    \vspace{1.5cm}
    {\color{myblue}\rule{\textwidth}{2pt}}\par
    \vspace{4cm}
    {\Large Ethan}\par
    \vspace{0.5cm}
    {\large 2026年11月}\par
    \vfill
\end{titlepage}

% ==================== 目录 ====================
\setcounter{page}{1}
\tableofcontents
\newpage

\vspace*{-0.5cm}
\begin{center}
    {\color{black}\heiti\Large\bfseries Four-State Method for Magnetic Interactions}
\end{center}
\vspace{0.3em}
\addcontentsline{toc}{chapter}{Four-State Method for Magnetic Interactions}
\markboth{Four-State Method for Magnetic Interactions}{}
\setcounter{section}{0}

\section{理论背景}

磁性体系的磁序及其稳定性由磁性位点之间的交换作用、自旋轨道耦合引起的各向异性等因素共同决定。
为了将第一性原理计算结果用于磁基态分析、自旋动力学或蒙特卡洛模拟，
需要将不同自旋构型的总能量映射到有效自旋哈密顿量，提取相应的磁相互作用参数。

在考虑各向同性对称交换、Dzyaloshinskii--Moriya（DM）相互作用和单轴单离子各向异性时，
可采用以下模型：
\[
\begin{aligned}
H={}&E_0+\sum_{\langle i,j\rangle}J_{ij}\mathbf S_i\cdot\mathbf S_j
+\sum_{\langle i,j\rangle}\mathbf D_{ij}\cdot(\mathbf S_i\times\mathbf S_j)+\sum_i A_i(\mathbf S_i\cdot\hat{\mathbf z}'_i)^2.
\end{aligned}
\]
其中，\(\langle i,j\rangle\) 表示每对磁性位点只计数一次，各项含义为：
\begin{itemize}
    \item \(E_0\)：与自旋方向无关的能量。
    \item \(J_{ij}\)：各向同性对称交换参数，决定两自旋平行或反平行排列的能量差。
    \item \(\mathbf D_{ij}\)：DM 相互作用矢量，描述反对称交换，可促使自旋发生倾斜并选择旋转手性。
    \item \(A_i\)：单轴单离子各向异性参数，描述自旋相对于局部轴 \(\hat{\mathbf z}'_i\) 的方向偏好。
\end{itemize}

本文采用上述哈密顿量的符号约定：\(J_{ij}>0\) 表示反铁磁交换，\(J_{ij}<0\) 表示铁磁交换；
对于 \(A_i(S_{iz'})^2\)，\(A_i<0\) 表示易轴，\(A_i>0\) 表示易平面。

各向同性交换 \(J\) 在不包含自旋轨道耦合（SOC）时即可存在。
DM 相互作用和本文讨论的单离子各向异性来源于 SOC，因此提取 \(D\) 和 \(A\) 时需要包含 SOC。

\section{四态法的基本原理}

四态法通过有选择地翻转目标自旋，将四种构型的总能量进行线性组合，
消去与目标参数无关的背景贡献。

对于两个目标位点，令其方向分别由 \(\sigma_1,\sigma_2=\pm1\) 表示。
在选定的自旋方向和上述模型下，四态能量可整理为
\[
E(\sigma_1,\sigma_2)=C+a\sigma_1+b\sigma_2+K\sigma_1\sigma_2.
\]
其中，\(C\) 为共同背景能量，\(a\)、\(b\) 包含目标自旋与固定背景的耦合贡献，
\(K\) 为待提取的双自旋项系数。
按照 \((+,+)\)、\((+,-)\)、\((-,+)\)、\((-,-)\) 的顺序编号，有
\[
E_1+E_4-E_2-E_3=4K.
\]
因此，目标自旋沿同一轴排列时，可提取 \(J_{12}\)；
目标自旋分别沿 \(x\)、\(y\) 排列时，在本文模型下可提取 \(D_{12}^{z}\)。

单离子各向异性的提取采用另一种配对方式：
先将目标自旋沿同一轴的正、反方向能量相加，
消去随目标自旋反向而变号的交换贡献，再比较两个轴方向的能量。
对于单位自旋和 \(H_A=A_1(S_{1z'})^2\)，有
\[
E_1+E_2=2C+2A_1,\qquad E_3+E_4=2C,
\]
从而得到
\[
A_1=\frac{E_1+E_2-E_3-E_4}{2}.
\]

\clearpage
\begingroup
\renewcommand{\arraystretch}{1.12}
\setlength{\abovedisplayskip}{5pt}
\setlength{\belowdisplayskip}{5pt}
\setlength{\abovedisplayshortskip}{3pt}
\setlength{\belowdisplayshortskip}{3pt}
\setlength{\topsep}{3pt}
\titlespacing*{\subsection}{0pt}{0.4em}{0.2em}
\noindent\begin{minipage}{\textwidth}
\section{四态自旋设置与计算公式}

采用单位自旋，\(|\mathbf S_i|=1\)。
注：“不变”表示该位点的自旋在同组四次计算中保持相同。

\noindent\begin{minipage}{\textwidth}
\subsection{对称交换参数 \(J_{12}\)}
\[
H_J=J_{12}\mathbf S_1\cdot\mathbf S_2.
\]
\begin{center}
\begin{tabular}{cccc}
\toprule
自旋态 & \(\mathbf S_1\) & \(\mathbf S_2\) & \(\Sother\) \\ \midrule
1 & \((0,0,1)\)  & \((0,0,1)\)  & 不变 \\
2 & \((0,0,1)\)  & \((0,0,-1)\) & 不变 \\
3 & \((0,0,-1)\) & \((0,0,1)\)  & 不变 \\
4 & \((0,0,-1)\) & \((0,0,-1)\) & 不变 \\
\bottomrule
\end{tabular}
\end{center}
\[
J_{12}=\frac{E_1+E_4-E_2-E_3}{4}.
\]
\end{minipage}
\par\medskip

\noindent\begin{minipage}{\textwidth}
\subsection{DM 参数 \(D_{12}^{z}\)}
\[
H_D=\mathbf D_{12}\cdot(\mathbf S_1\times\mathbf S_2).
\]
\begin{center}
\begin{tabular}{cccc}
\toprule
自旋态 & \(\mathbf S_1\) & \(\mathbf S_2\) & \(\Sother\) \\ \midrule
1 & \((1,0,0)\)  & \((0,1,0)\)  & 不变 \\
2 & \((1,0,0)\)  & \((0,-1,0)\) & 不变 \\
3 & \((-1,0,0)\) & \((0,1,0)\)  & 不变 \\
4 & \((-1,0,0)\) & \((0,-1,0)\) & 不变 \\
\bottomrule
\end{tabular}
\end{center}
论文采用沿 \(z\) 轴的固定背景，计算须包含 SOC。
\[
D_{12}^{z}=\frac{E_1+E_4-E_2-E_3}{4}.
\]
\end{minipage}
\par\medskip

\noindent\begin{minipage}{\textwidth}
\subsection{单离子各向异性参数 \(A_1\)}
以下分量采用局部坐标系 \(x',y',z'\)：
\[
H_A=A_1(S_{1z'})^2.
\]
\begin{center}
\begin{tabular}{cccc}
\toprule
自旋态 & \(\mathbf S_1\) & \(\mathbf S_2\) & \(\Sother\) \\ \midrule
1 & \((0,0,1)\)  & 不变 & 不变 \\
2 & \((0,0,-1)\) & 不变 & 不变 \\
3 & \((1,0,0)\)  & 不变 & 不变 \\
4 & \((-1,0,0)\) & 不变 & 不变 \\
\bottomrule
\end{tabular}
\end{center}
论文采用沿 \(y'\) 的固定背景，计算须包含 SOC。
\[
A_1=\frac{E_1+E_2-E_3-E_4}{2}.
\]
\end{minipage}


\end{minipage}
\endgroup
\clearpage
\section{计算条件与适用范围}

四态能量映射要求同一组计算对应同一个有效自旋模型，具体应保持：
\begin{itemize}
    \item \textbf{原子结构一致：}采用相同晶胞和原子坐标，不能将四态分别弛豫后的能量直接用于提取同一参考结构的参数。
    \item \textbf{计算设置一致：}采用相同的交换关联泛函、赝势、截断能、\(k\) 点、DFT+\(U\) 参数及收敛标准。
    \item \textbf{背景自旋一致：}同组四态中，未被指定改变的各位点自旋均保持原方向；不同参数计算组之间不要求背景完全相同。
    \item \textbf{目标方向受控：}应检查实际收敛后的磁矩方向，必要时使用约束自旋计算，防止目标构型自行转变。
    \item \textbf{局域磁矩基本稳定：}四态法采用固定自旋长度模型；若不同构型的局域磁矩大小显著变化，映射结果需要进一步检查。
    \item \textbf{超胞足够大：}避免目标位点与周期镜像之间的作用混入所提取的单条磁性键参数。
\end{itemize}

以上公式针对所列哈密顿量。
若体系存在显著的对称各向异性交换、高阶自旋作用等额外项，
需要扩展模型和自旋构型，检查这些项对参数提取的影响。
采用单位自旋时，所得参数已包含自旋长度归一化的影响，
与采用 \(|\mathbf S_i|=S\) 定义的参数比较时应统一约定。

\vspace{0.5cm}

\subsection*{参考文献}
\addcontentsline{toc}{subsection}{参考文献}

H. J. Xiang, E. J. Kan, S.-H. Wei, M.-H. Whangbo, and X. G. Gong,
``Predicting the spin-lattice order of frustrated systems from first principles,''
\textit{Physical Review B} \textbf{84}, 224429 (2011).
\par\noindent
\href{https://doi.org/10.1103/PhysRevB.84.224429}
{DOI: 10.1103/PhysRevB.84.224429.}

\end{document}